\documentclass[aspectratio=169]{beamer}
\usepackage[utf8]{inputenc}
\usepackage[T1]{fontenc}
\usepackage[spanish]{babel}
\usepackage{graphicx}
\usepackage{tikz}
\usepackage{listings}
\usepackage{xcolor}

% Tema
\usetheme{Madrid}
\usecolortheme{default}

% Colores personalizados
\definecolor{unsazul}{RGB}{30,60,120}
\definecolor{unsverde}{RGB}{0,100,0}

\setbeamercolor{structure}{fg=unsazul}
\setbeamercolor{palette primary}{bg=unsazul,fg=white}
\setbeamercolor{palette secondary}{bg=unsazul!80,fg=white}
\setbeamercolor{palette tertiary}{bg=unsazul!60,fg=white}

% Configuración de listings (código)
\lstset{
    basicstyle=\ttfamily\small,
    keywordstyle=\color{unsazul}\bfseries,
    commentstyle=\color{gray}\itshape,
    stringstyle=\color{unsverde},
    showstringspaces=false,
    breaklines=true,
    frame=single,
    backgroundcolor=\color{gray!10}
}

% Información del documento
\title{Métodos de Análisis de Datos 1}
\subtitle{Clase 2: Notebooks y Buenas Prácticas}
\author{Universidad Nacional del Sur}
\date{}

\begin{document}

% Portada
\begin{frame}
\titlepage
\end{frame}

% Contenidos
\begin{frame}{Contenidos de hoy}
\tableofcontents
\end{frame}

\section{Jupyter Notebooks}

\begin{frame}{¿Qué son los Jupyter Notebooks?}
\begin{block}{Definición}
Documentos \textbf{interactivos} que combinan:
\begin{itemize}
    \item Código ejecutable (Python, R, Julia...)
    \item Texto narrativo (Markdown, LaTeX)
    \item Visualizaciones (gráficos, tablas)
    \item Resultados (outputs)
\end{itemize}
\end{block}

\pause

\begin{exampleblock}{¿Por qué notebooks?}
\begin{itemize}
    \item Facilitan la exploración interactiva
    \item Documentan el proceso de análisis
    \item Permiten compartir resultados reproducibles
\end{itemize}
\end{exampleblock}
\end{frame}

\begin{frame}{Estructura de un notebook}
\begin{block}{Celdas de código}
\begin{itemize}
    \item Ejecutan código (Python, R, etc.)
    \item Muestran resultados inmediatamente
    \item Pueden generar gráficos y tablas
\end{itemize}
\end{block}

\begin{block}{Celdas Markdown}
\begin{itemize}
    \item Texto formateado (negritas, cursivas, listas)
    \item Títulos y secciones
    \item Ecuaciones matemáticas con \LaTeX
    \item Imágenes y links
\end{itemize}
\end{block}
\end{frame}

\begin{frame}{Atajos de teclado útiles}
\begin{block}{Modo comando (Esc)}
\begin{itemize}
    \item \texttt{A}: Insertar celda arriba
    \item \texttt{B}: Insertar celda abajo
    \item \texttt{D D}: Borrar celda
    \item \texttt{M}: Convertir a Markdown
    \item \texttt{Y}: Convertir a código
\end{itemize}
\end{block}

\begin{block}{Modo edición (Enter)}
\begin{itemize}
    \item \texttt{Shift + Enter}: Ejecutar y avanzar
    \item \texttt{Ctrl + Enter}: Ejecutar sin avanzar
    \item \texttt{Tab}: Autocompletar
\end{itemize}
\end{block}
\end{frame}

\section{Sintaxis básica de Python}

\begin{frame}[fragile]{Python: Variables y tipos de datos}
\begin{lstlisting}[language=Python]
# Variables y tipos basicos
x = 10              # int
y = 3.14            # float
nombre = "Ana"      # str
activo = True       # bool

# Listas
numeros = [1, 2, 3, 4, 5]
mezclado = [1, "dos", 3.0, True]

# Diccionarios
persona = {
    "nombre": "Ana",
    "edad": 25,
    "ciudad": "Bahia Blanca"
}
\end{lstlisting}
\end{frame}

\begin{frame}[fragile]{Python: Control de flujo}
\begin{lstlisting}[language=Python]
# Condicionales
if x > 5:
    print("Mayor a 5")
elif x == 5:
    print("Igual a 5")
else:
    print("Menor a 5")

# Bucles
for i in range(5):
    print(i)

# List comprehension
cuadrados = [x**2 for x in range(10)]
\end{lstlisting}
\end{frame}

\begin{frame}[fragile]{Python: Funciones}
\begin{lstlisting}[language=Python]
# Funcion simple
def saludar(nombre):
    return f"Hola, {nombre}!"

# Funcion con argumentos por defecto
def potencia(base, exponente=2):
    return base ** exponente

# Lambda (funciones anonimas)
cuadrado = lambda x: x**2

# Ejemplo de uso
resultado = potencia(3)      # 9
resultado = potencia(3, 3)   # 27
\end{lstlisting}
\end{frame}

\section{Sintaxis básica de R}

\begin{frame}[fragile]{R: Variables y estructuras}
\begin{lstlisting}[language=R]
# Variables
x <- 10
y <- 3.14
nombre <- "Ana"

# Vectores
numeros <- c(1, 2, 3, 4, 5)
secuencia <- 1:10

# Listas
persona <- list(
    nombre = "Ana",
    edad = 25,
    ciudad = "Bahia Blanca"
)

# Data frames
df <- data.frame(
    x = 1:5,
    y = c(2, 4, 6, 8, 10)
)
\end{lstlisting}
\end{frame}

\begin{frame}[fragile]{R: Funciones}
\begin{lstlisting}[language=R]
# Funcion simple
saludar <- function(nombre) {
    paste("Hola,", nombre, "!")
}

# Funcion con argumentos por defecto
potencia <- function(base, exponente = 2) {
    base ^ exponente
}

# Funcion apply (vectorizada)
cuadrados <- sapply(1:10, function(x) x^2)

# Ejemplo de uso
resultado <- potencia(3)      # 9
resultado <- potencia(3, 3)   # 27
\end{lstlisting}
\end{frame}

\section{Buenas prácticas}

\begin{frame}{Buenas prácticas de programación}
\begin{block}{Código limpio}
\begin{itemize}
    \item Nombres descriptivos de variables
    \item Comentarios claros y concisos
    \item Indentación consistente
    \item Funciones pequeñas y enfocadas
\end{itemize}
\end{block}

\begin{block}{Reproducibilidad}
\begin{itemize}
    \item Documentar dependencias y versiones
    \item Usar semillas aleatorias (\texttt{random.seed()}, \texttt{set.seed()})
    \item Guardar outputs intermedios
    \item Registrar parámetros de experimentos
\end{itemize}
\end{block}
\end{frame}

\begin{frame}{Organización de proyectos}
\begin{block}{Estructura recomendada}
\begin{itemize}
    \item \texttt{data/}: Datos crudos y procesados
    \item \texttt{notebooks/}: Jupyter notebooks
    \item \texttt{src/}: Código fuente (módulos y funciones)
    \item \texttt{outputs/}: Resultados, gráficos, reportes
    \item \texttt{README.md}: Descripción del proyecto
    \item \texttt{requirements.txt} o \texttt{environment.yml}
\end{itemize}
\end{block}

\begin{alertblock}{Importante}
\textbf{Nunca} subir datos sensibles o credenciales a GitHub
\end{alertblock}
\end{frame}

\section{Git y GitHub}

\begin{frame}[fragile]{Git: Control de versiones}
\begin{block}{Comandos básicos}
\begin{lstlisting}[language=bash]
# Inicializar repositorio
git init

# Agregar archivos al staging
git add archivo.py
git add .

# Commit (guardar cambios)
git commit -m "Mensaje descriptivo"

# Ver estado
git status

# Ver historial
git log
\end{lstlisting}
\end{block}
\end{frame}

\begin{frame}[fragile]{GitHub: Colaboración}
\begin{block}{Comandos para trabajo remoto}
\begin{lstlisting}[language=bash]
# Clonar repositorio
git clone https://github.com/usuario/repo.git

# Subir cambios
git push origin main

# Bajar cambios
git pull origin main

# Crear rama
git branch nueva-funcionalidad
git checkout nueva-funcionalidad
\end{lstlisting}
\end{block}

\begin{exampleblock}{Flujo de trabajo típico}
\texttt{git pull} → modificar → \texttt{git add} → \texttt{git commit} → \texttt{git push}
\end{exampleblock}
\end{frame}

\section{R Markdown}

\begin{frame}[fragile]{R Markdown: Reportes reproducibles}
\small
\begin{block}{Estructura}
\begin{lstlisting}
---
title: "Mi Analisis"
author: "Tu Nombre"
date: "2026-03-16"
output: html_document
---
# Introduccion
Este es texto normal en Markdown.
    ```{r}
# Este es codigo R
x <- 1:10
mean(x)
```
El promedio es `r mean(x)`.
\end{lstlisting}
\end{block}

\small
El código entre \texttt{```\{r\}} se ejecuta al generar el reporte
\end{frame}

\section{Caso de estudio}

\begin{frame}{Caso: Dataset Iris}
\begin{columns}[t]
\column{0.5\textwidth}
\begin{block}{Descripción}
\begin{itemize}
    \item 150 flores de iris
    \item 3 especies diferentes
    \item 4 medidas por flor:
    \begin{itemize}
        \item Largo de sépalo
        \item Ancho de sépalo
        \item Largo de pétalo
        \item Ancho de pétalo
    \end{itemize}
\end{itemize}
\end{block}

\column{0.5\textwidth}
\begin{block}{Objetivo}
Recorrer todo el ciclo:
\begin{enumerate}
    \item Cargar datos
    \item Explorar
    \item Visualizar
    \item Modelo simple
    \item Interpretar
\end{enumerate}
\end{block}
\end{columns}

\vspace{1em}

\begin{alertblock}{En el notebook de práctica}
Veremos el caso completo implementado en Python y R
\end{alertblock}
\end{frame}

\section{Resumen}

\begin{frame}{Resumen de la clase}
\begin{itemize}
    \item \textbf{Jupyter notebooks}: combinan código, texto y visualizaciones
    \item \textbf{Python y R}: sintaxis básica, variables, funciones
    \item \textbf{Buenas prácticas}: código limpio, reproducibilidad, organización
    \item \textbf{Git y GitHub}: control de versiones, colaboración
    \item \textbf{R Markdown}: reportes reproducibles en R
\end{itemize}

\vspace{1em}

\begin{block}{Próxima clase (Unidad 1)}
\begin{itemize}
    \item Fuentes de datos
    \item Carga de datos en Python y R
    \item Exploración inicial
\end{itemize}
\end{block}
\end{frame}

\begin{frame}{Tarea para la próxima clase}
\begin{enumerate}
    \item Completar instalación de herramientas
    \item Crear cuenta en GitHub y primer repositorio
    \item Configurar Git localmente
    \item Hacer el notebook de práctica de U0
    \item Subir el notebook al repositorio
\end{enumerate}

\vspace{1em}

\begin{alertblock}{Importante}
El notebook de práctica está disponible en el campus virtual
\end{alertblock}
\end{frame}


\end{document}
